\section{Toward a Blockchain-Native Datastore}
\label{sec:chain}

The optimizations of Sections~\ref{sec:build}--\ref{sec:eval} make Datastore a
better single-node engine. This section sketches a more ambitious trajectory:
operating Datastore as a self-hosted service that \emph{extends S-Chain} --- the
network's decentralized object-storage substrate --- co-located on the very same
nodes that serve S-Chain, with storage and compute separated over that shared
substrate. We are deliberate about what this is \emph{not}: Datastore is not a
new public chain with its own letter and validator economy (the network's chain
letters are otherwise assigned, e.g.\ the DEX). It is a query/compute engine
layered on S-Chain, using S-Chain for durable storage and the nodes' existing
Quasar consensus for coordination --- a permissioned, operator-run (or
LLM-consensus-node-run) service, not a permissionless chain. We report a working
proof-of-concept for the foundational layer and the validated integration path
for the rest.

\subsection{Separation of storage and compute}

The dominant cloud-analytics architecture (Snowflake, Databricks, ClickHouse
Cloud) separates a durable object-storage tier from stateless, elastic compute.
ClickHouse's \texttt{ReplicatedMergeTree} already factors cleanly along this
line: a small ordered \emph{replication log} and large, content-addressed
\emph{data parts}. We exploit this by pointing Datastore's MergeTree storage
policy at an S3-compatible object store --- here SeaweedFS, the storage engine
behind the network's ``S-Chain.''

\paragraph{Proof of concept.} We loaded a $5$M-row MergeTree table whose
storage policy targets SeaweedFS over S3. Verified outcome: all five parts
report disk \texttt{seaweed}; the local data directory holds $24$\,KB of
metadata; SeaweedFS holds $101$ content-addressed objects; and queries execute
correctly against the object-backed table (aggregation $25$\,ms, filter
$7$\,ms over $5$M rows). Resident memory was $68$\,MB versus $130$\,MB for an
upstream binary on local disk.

\paragraph{Honest cost.} Object storage loses a warm-cache drag race to local
disk by $2$--$4\times$ on single queries --- a round trip versus a page-cache
hit. The win is not raw warm latency; it is (i) a stateless compute tier
($24$\,KB local vs.\ $99$\,MB), enabling $N$ disposable compute nodes over one
copy of data --- which a local-disk engine structurally cannot do --- and (ii)
roughly half the resident memory. A local filesystem cache over the object disk
(the ClickHouse-Cloud technique) recovers most of the warm-latency gap.

\paragraph{Elastic scale-out, verified.} The stateless-compute claim is the one
a local-disk engine cannot make, so we verified it directly. With a
\texttt{ReplicatedMergeTree} over the S-Chain object store and zero-copy
replication enabled, we ran two independent compute nodes coordinated by Keeper:
node~1 inserted $5$M rows; node~2, started afterward, served the identical
$5{,}000{,}000$ rows. Critically, the object count in S-Chain rose from $101$ to
$102$ --- a single manifest object, \emph{not} a second copy of the data ---
confirming that the two replicas share one physical copy via zero-copy. Two
compute nodes, one copy of data on decentralized storage: the elastic-compute
property, demonstrated. (The single non-obvious requirement was configuring the
inter-node part-transfer endpoint; see Section~\ref{sec:zapnative}.)

\subsection{Replacing Raft with leaderless post-quantum consensus}

Coordination of the replication log is performed today by an embedded
Keeper on NuRaft (Raft): crash-fault-tolerant, leader-based, classical
cryptography. Replacing the Raft backend with the network's leaderless
consensus (Quasar) yields three properties at once: leaderless resilience (no
single writer to fail over), Byzantine fault tolerance over a bonded validator
set, and --- because Quasar finalizes with ML-DSA --- \emph{post-quantum
finality}. This is not speculative: a sibling system already runs SQLite under
the same leaderless Quasar coordination (writer-pin, all nodes equal
validators), establishing the pattern. The realistic integration is a
coordination service that speaks the ZooKeeper API Datastore already expects but
backs the ordered log with Quasar, so the engine itself is unchanged. The result
is, to our knowledge, the first \emph{post-quantum-native} columnar OLAP engine.

\subsection{Extending S-Chain, not minting a new chain}

Co-locating Datastore with S-Chain rather than standing up a separate chain is
the simpler and more honest design: bulk columnar data is far too large to pass
through consensus (the network caps a gossiped block at $2$\,MB), so a
``datastore chain'' would in any case put \emph{only} small commitments on-chain
and keep the data off-chain --- which is exactly what running beside S-Chain
already gives us, without a new validator set or chain letter. The division of
labor is therefore: the columnar \emph{parts} are content-addressed objects in
S-Chain storage; the small \emph{commitments} (part-manifest roots, DDL,
schema/version pointers) are ordered by the nodes' existing Quasar consensus; and
reads are served by the co-located Datastore compute engine. Two on-disk
precedents make the engine integration concrete without inventing anything: a
sibling component already binds a C++ execution engine behind a thin Go shim via
FFI (``valid as long as state roots match''), and another already serves a
read-only query surface over a handler while keeping bulk state out of blocks.
Datastore reuses both shapes as a \emph{service on S-Chain nodes} --- the C++
engine for execution, a query surface for reads --- with ZAP
(Section~\ref{sec:arch}) as the wire throughout. Should a workload genuinely want
a public, permissionless analytics chain later (an LLM-consensus quorum
publishing verifiable aggregates, say), the same commitments-on-chain /
parts-off-chain split promotes to a first-class chain unchanged; but that is an
option, not the default.

\subsection{Permissionless operation with confidential compute}

Open, untrusted query execution is the residual hard problem: a serving node's
\texttt{SELECT} result is not verifiable by re-execution for arbitrary SQL, and
ZK-SQL is not practical at OLAP scale. The pragmatic resolution is
attestation rather than proof: run compute inside a hardware trusted execution
environment (e.g.\ confidential-compute GPUs), so a remote party verifies an
attestation of the enclave rather than re-deriving the query. Combined with
content-addressed parts (integrity) and Quasar-ordered manifests (a verifiable
write history), this yields a system that is open to read and use, bonded to
write and serve, and private under attestation --- without requiring verifiable
arbitrary-query mathematics.

\subsection{One transport: ZAP-native, end to end}
\label{sec:zapnative}

A conventional analytics deployment carries three distinct wire protocols: a
client protocol, a gRPC service surface, an inter-node replication transport
(HTTP), and a ZooKeeper wire for coordination. Our objective is to collapse all
of these onto one substrate --- ZAP, the Cap'n-Proto-derived transport with
zero-copy framing and promise pipelining (Section~\ref{sec:arch}). The current
state and the remaining work:

\begin{table}[H]
\centering
\small
\begin{tabular}{@{}lll@{}}
\toprule
Wire & Today & ZAP-native target \\
\midrule
gRPC service     & \textbf{removed}      & --- (gone) \\
Inter-node parts & HTTP (DataPartsExch.) & ZAP \\
Coordination log & ZooKeeper / Keeper    & Quasar (ZAP) \\
Client query     & native TCP / ZAP stub & SQL-over-ZAP \\
\bottomrule
\end{tabular}
\caption{Collapsing four wires onto ZAP. gRPC is already gone; the inter-node
HTTP part-transfer and the ZooKeeper coordination wire are the remaining
targets.}
\label{tab:zapnative}
\end{table}

The gRPC surface and ArrowFlight are already deleted: the binary contains no
\texttt{GRPCServer} and no \texttt{USE\_GRPC} code paths. The inter-node
part-transfer --- the \texttt{InterserverIOHTTPHandler}/\texttt{DataPartsExchange}
that the scale-out result above relies on --- runs over HTTP today and is the
first migration to ZAP, since part fetch is exactly the kind of large,
streamable, pipelineable transfer ZAP's opaque-payload model was designed for.
The coordination wire (the $17$-file ZooKeeper client talking to Keeper) is the
second: replaced by the Quasar-backed coordination service of
Section~\ref{sec:chain}, ``ZooKeeper'' disappears as a wire while the engine's
expectations are preserved. When both land, alongside the SQL-over-ZAP query
path, the platform is one transport deep --- no HTTP, no gRPC, no ZooKeeper ---
which is also precisely the precondition for the unified two-engine platform of
Section~\ref{sec:twoengine}.

\paragraph{Phasing.} The product engine ships today (Sections~\ref{sec:build}--%
\ref{sec:eval}); storage/compute separation \emph{and elastic scale-out with
zero-copy} over S-Chain are demonstrated above; the Raft$\rightarrow$Quasar
coordination swap, the inter-node-HTTP$\rightarrow$ZAP migration, and the chain
VM are the near-term build, each with an on-disk precedent; verifiable/attested
permissionless reads are the research frontier. Each phase preserves
compatibility and pays only for what the workload uses.
